\section{Networking and addressing: \FDM, \PDNS, mesh resolution, and private deployments}
\label{sec:networking}

An operator can point a browser or an API client at an application running somewhere on the
several-thousand-node \Flux fleet and get an answer. That works today, and the path by which it
works is worth naming precisely, because every hop on it is presently operated by the \Flux
Foundation. This is not a claim about node operators: the compute itself is dispersed across
independent Cumulus, Nimbus and Stratus operators exactly as advertised. It is a claim about
\emph{reachability} — DNS resolution, TLS termination, and the reverse proxy that decides which
container answers a given hostname all run on Foundation-controlled infrastructure, discovered
through a single Foundation-operated API with no on-chain or peer-to-peer fallback. Naming this
path precisely, hop by hop, is what makes the decentralization roadmap in \S\ref{sec:governance}
and elsewhere legible: a reader cannot evaluate a plan to decentralize a system without first
knowing exactly what is centralized in it today.

\subsection{The reachability path, end to end}
\label{sec:networking-path}

Application placement itself is decided elsewhere, by \FluxOS (\S\ref{sec:fluxos}); nothing in
this section governs \emph{where} an application runs. What this section governs is how a client
finds the IP:port that \FluxOS placed it on. \shipped

\begin{enumerate}
  \item \textbf{Discovery.} \FDM (\texttt{flux-\allowbreak{}domain-\allowbreak{}manager}) does not talk to \fluxd or to any
  peer-to-peer layer to learn what applications exist or where they run. It polls a single
  Flux-Foundation-operated REST API, \texttt{https:/\allowbreak{}/\allowbreak{}api.\allowbreak{}runonflux.\allowbreak{}io}, for
  \texttt{apps/\allowbreak{}globalappsspecifications} and \texttt{apps/locations}
  \srccode{flux-domain-manager/\allowbreak{}src/\allowbreak{}services/\allowbreak{}flux/\allowbreak{}dataFetcher.js}{47-71}. This is FDM's only
  source for "what apps exist and where"; no fallback source exists in the codebase
  \srccode{flux-domain-manager/\allowbreak{}src/\allowbreak{}services/\allowbreak{}flux/\allowbreak{}dataFetcher.js}{47-91}.
  \item \textbf{Polling cadence.} Three independent, self-rescheduling timers run per FDM process:
  \texttt{globalAppSpecs} at a default 30{,}000\,ms (conditional on HTTP \texttt{ETag}), \texttt{
  apps\allowbreak{}Locations} hardcoded to 10{,}000\,ms regardless of change, and \texttt{permMessages} at a
  default 120{,}000\,ms \srccode{flux-domain-manager/\allowbreak{}src/\allowbreak{}services/\allowbreak{}flux/\allowbreak{}dataFetcher.js}{47-91,409-417,689-721}.
  \item \textbf{Naming.} For a native application $\app$ with port $P$, FDM computes
  \texttt{$\app$\_$P$.app2.runonflux.io} (per port) and \texttt{$\app$.app2.runonflux.io} (main
  alias) by string concatenation — \texttt{app.} rather than \texttt{app2.} on the production
  shard \srccode{flux-domain-manager/\allowbreak{}src/\allowbreak{}services/\allowbreak{}domain/\allowbreak{}index.js}{9-32}. An operator may also
  set a custom domain in the application spec, which FDM verifies against its own IP via
  \texttt{dns.resolve4} before requesting a certificate for it
  \srccode{flux-domain-manager/\allowbreak{}src/\allowbreak{}services/\allowbreak{}domain/\allowbreak{}cert.js}{128-157}.
  \item \textbf{DNS.} A hostname is served either by Cloudflare, used strictly as authoritative
  "DNS-only" (grey-cloud) DNS — FDM actively deletes any record it finds with \texttt{proxied =\allowbreak{}=\allowbreak{}=\allowbreak{}
  true} — or by the self-hosted \PDNS zone service, selected per FDM host by
  \texttt{config.\allowbreak{}cloudflare.\allowbreak{}enabled}\slash\texttt{config.\allowbreak{}p\allowbreak{}DNS.\allowbreak{}enabled}
  \srccode{flux-domain-manager/\allowbreak{}src/\allowbreak{}services/\allowbreak{}domain/\allowbreak{}cert.js}{170-181}
  \srccode{flux-domain-manager/\allowbreak{}config/\allowbreak{}default.js}{32-47}. The evidence surveyed establishes how
  each backend behaves but not, from code alone, which one is authoritative for
  \texttt{runonflux.io} at any given moment: the checked-in default config is itself a
  certificate-only variant with \texttt{cloudflare.enabled:\ true}, and per-host Ansible
  templating can override this per fleet member \srccode{flux-domain-manager/\allowbreak{}config/\allowbreak{}default.js}{38,46}.
  \item \textbf{TLS termination and routing.} The FDM host's own \texttt{HAProxy} instance
  terminates TLS, binding \texttt{*:443} against every \texttt{.pem} file under
  \texttt{/\allowbreak{}etc/\allowbreak{}ssl/\allowbreak{}fluxapps/\allowbreak{}} \srccode{flux-domain-manager/\allowbreak{}src/\allowbreak{}services/\allowbreak{}haproxyTemplate.js}{81-134},
  then routes on the \texttt{Host} header via an ACL to the matching backend, forwarding to one of
  the application's container IPs at its declared port
  \srccode{flux-domain-manager/\allowbreak{}src/\allowbreak{}services/\allowbreak{}haproxyTemplate.js}{260-319}. Per-server health
  checks (\S\ref{sec:networking-runloop}) decide which of an application's IPs are actually in that
  backend at any moment.
  \item \textbf{Game-application shortcut.} For roughly two dozen "G-mode" game-server application
  types (Minecraft, Rust, Valheim, and similar), a separate, single-instance service,
  \texttt{flux-\allowbreak{}apps-\allowbreak{}dns-\allowbreak{}manager}, reads the master IP FDM already elected into its own
  \texttt{/appips} output and publishes it directly as an \texttt{A} record via the DNS Gateway
  (\S\ref{sec:networking-gateway}) — client traffic, often UDP, goes node-to-node and never passes
  through HAProxy \srccode{flux-apps-dns-manager/\allowbreak{}src/\allowbreak{}services/\allowbreak{}appsDnsManager.js}{202-255}
  \srcdoc{flux-apps-dns-manager/README.md:3-15}.
\end{enumerate}

\begin{figure}[H]
  \centering
  \begin{tikzpicture}[
    font=\footnotesize,
    node distance=6mm,
    box/.style={draw, rounded corners=2pt, minimum height=9mm, align=center,
                inner xsep=3pt, text width=23mm},
    flux/.style={box, fill=black!4},
    found/.style={box, fill=black!12, very thick},
    lbl/.style={font=\scriptsize, align=center, inner sep=1pt},
    arr/.style={-{Latex[length=2mm]}, thick},
  ]
    \node[box]   (client) {client};
    \node[found, right=10mm of client] (dns) {DNS\\[-1pt]{\scriptsize Cloudflare \emph{or} \PDNS}};
    \node[found, right=10mm of dns]    (ha)  {HAProxy\\[-1pt]{\scriptsize on the elected \FDM host}};
    \node[flux,  right=10mm of ha]     (app) {application\\[-1pt]{\scriptsize container IP:port}};

    \draw[arr] (client) -- node[lbl, above] {name} (dns);
    \draw[arr] (dns) -- node[lbl, above] {\texttt{A} record} (ha);
    \draw[arr] (ha)  -- node[lbl, above] {proxied} (app);

    % HAProxy-bypassing shortcut for game applications
    \draw[arr, dashed] (dns.south) .. controls +(0,-11mm) and +(0,-11mm) ..
      node[lbl, below, pos=0.5] {\texttt{flux-\allowbreak{}apps-\allowbreak{}dns-\allowbreak{}manager}: elected master IP as a direct
        \texttt{A} record\\[-1pt](game applications; bypasses HAProxy)} (app.south);

    %% Named and included in the fit: as a bare label above `ha' it printed on top of
    %% the dashed box's own "Foundation-operated" label, which sits at the same height.
    \node[lbl, above=1.5mm of ha] (tls) {TLS termination, Host-header ACL,\\[-1pt]per-server health check};
    \node[draw, dashed, rounded corners=2pt, inner sep=3mm, fit=(dns)(ha)(tls),
          label={[lbl, yshift=1.5mm]above:{Foundation-operated}}] {};
  \end{tikzpicture}
  \caption{The reachability path from client to a running application, and the game-application
  shortcut that bypasses \FDM's HAProxy entirely. Both intermediate hops are Foundation-operated
  \shipped; \S\ref{sec:networking-inventory} inventories what each depends on and what happens
  when it is unavailable.}
  \label{fig:reachability-path}
\end{figure}

The \texttt{*.\allowbreak{}app.\allowbreak{}runonflux.\allowbreak{}io}\slash\texttt{*.\allowbreak{}app2.\allowbreak{}runonflux.\allowbreak{}io} wildcard certificate that step 5
depends on is treated by \FDM's own code as always present — the presence check short-circuits to
\texttt{true} for any domain ending in these suffixes, so \FDM never calls \texttt{certbot} for
them \srccode{flux-domain-manager/\allowbreak{}src/\allowbreak{}services/\allowbreak{}domain/\allowbreak{}cert.js}{18-22}.
\decided{the wildcard certificate is managed by InFlux Technologies USA LLC and renewed
periodically against the \texttt{runonflux.io} domain \srcintent{08-answers-round-2.md, rounds
24--25}. This is less of a concentration than it first reads. Behind any certificate there must be
some company managing the domain --- that is a property of the web PKI rather than of \Flux{} --- and
what matters is whether the arrangement can be replicated. It can: \texttt{flux-domain-manager} is
public under AGPL-3.0, the application specification already admits tenant-set custom domains
\srccode{flux/\allowbreak{}ZelBack/\allowbreak{}src/\allowbreak{}services/\allowbreak{}appsService.js}{passim}, and any operator may run
their own \FDM{}, their own \PDNS{} and serve applications from a domain they control.

Discovery is replicable too, and this paper's earlier reading of
\texttt{api.\allowbreak{}runonflux.\allowbreak{}io} conflated two different claims. The endpoint is not a
Foundation-only service: \texttt{/apps/\allowbreak{}globalappsspecifications} is an ordinary \FluxOS{}
route served by \emph{every} node \shipped~\srccode{flux/\allowbreak{}ZelBack/\allowbreak{}src/\allowbreak{}routes.js}{402}, and
\FDM{} takes its host as a constructor parameter, \texttt{flux\allowbreak{}Api\allowbreak{}BaseUrl}, whose
shipped value merely happens to be that name
\shipped~\srccode{flux-domain-manager/\allowbreak{}src/\allowbreak{}services/\allowbreak{}flux/\allowbreak{}dataFetcher.js}{886}. Pointing an
independent \FDM{} at any of the \num{6489} nodes, or at a bare IP, is a configuration change rather
than a fork --- and the code is AGPL-3.0 if a fork is wanted anyway.

What survives is narrower and is a resilience property, not a centralisation one: as shipped there is
\emph{one} base URL and \emph{no automatic failover}, so an \FDM{} whose configured endpoint becomes
unreachable stalls on its last-known data rather than trying another node. That is worth fixing and
is recorded in \Cref{tab:centralization} on those terms. \textbf{What is centralised here is the
default, not the capability.}}

\subsection{\FDM's run loop}
\label{sec:networking-runloop}

\FDM's inputs are the three polled feeds above (application specs, locations, permanent
messages); its output is a generated \texttt{HAProxy} configuration file, validated with
\texttt{haproxy -\allowbreak{}f .\allowbreak{}.\allowbreak{}.\allowbreak{} -\allowbreak{}c} before it is ever installed — a failed validation leaves the
previously-loaded configuration serving, not a hard outage
\srccode{flux-domain-manager/\allowbreak{}src/\allowbreak{}services/\allowbreak{}domainService.js}{196-200,355-359}
\srccode{flux-domain-manager/\allowbreak{}src/\allowbreak{}services/\allowbreak{}haproxyTemplate.js}{626-646}. Reconfiguration of the
main HAProxy pool (the \texttt{home}\slash\texttt{api}\slash\texttt{cloud} UI/API backends) is driven purely
by a 30{,}000\,ms self-rescheduling timer, independent of any external event
\srccode{flux-domain-manager/\allowbreak{}src/\allowbreak{}services/\allowbreak{}domainService.js}{96-210}; that regeneration throws and
aborts the cycle if fewer than 10 Stratus candidates pass its 8-stage health chain, keeping the
old configuration running rather than publishing an under-populated pool
\srccode{flux-domain-manager/\allowbreak{}src/\allowbreak{}services/\allowbreak{}domainService.js}{110-183}. Reconfiguration of
per-application backends is driven by the \texttt{apps\allowbreak{}Locations\allowbreak{}Updated} event fired on every
10-second location poll, which re-checks every location IP with a per-application-type coded
health check and writes into the backend only the IPs that pass
\srccode{flux-domain-manager/\allowbreak{}src/\allowbreak{}services/\allowbreak{}domainService.js}{730-767}. \shipped

\begin{table}[H]
\centering
\small
\renewcommand{\arraystretch}{1.15}
\caption{\FDM/\PDNS reachability-path health-check and reconfiguration constants.}
\label{tab:fdm-health-constants}
\begin{tabular}{lll}
\toprule
Mechanism & Constants & Source \\
\midrule
Generic app backend, HAProxy-level & \texttt{inter 3s fall 2 rise 2 fastinter 500} & \srccode{flux-domain-manager/\allowbreak{}src/\allowbreak{}services/\allowbreak{}haproxyTemplate.js}{311} \\
\quad detection floor & $\approx$6\,s down (2$\times$3\,s) & \srccode{flux-domain-manager/\allowbreak{}src/\allowbreak{}services/\allowbreak{}haproxyTemplate.js}{311} \\
Main UI/API backend, HAProxy-level & \texttt{inter 10s fall 2 rise 3 maxconn 100} & \srccode{flux-domain-manager/\allowbreak{}src/\allowbreak{}services/\allowbreak{}haproxyTemplate.js}{384,408,427} \\
\quad detection floor & $\approx$20\,s down, 30\,s up & \srccode{flux-domain-manager/\allowbreak{}src/\allowbreak{}services/\allowbreak{}haproxyTemplate.js}{384,408,427} \\
G-mode sticky-master retries & 3 attempts, 3{,}000\,ms apart & \srccode{flux-domain-manager/\allowbreak{}src/\allowbreak{}services/\allowbreak{}domainService.js}{29-30} \\
G-mode sticky-unhealthy grace & 90{,}000\,ms since last healthy & \srccode{flux-domain-manager/\allowbreak{}src/\allowbreak{}services/\allowbreak{}domainService.js}{31} \\
\texttt{appsLocations} poll (control plane) & 10{,}000\,ms, hardcoded & \srccode{flux-domain-manager/\allowbreak{}src/\allowbreak{}services/\allowbreak{}flux/\allowbreak{}dataFetcher.js}{66,707} \\
Main HAProxy pool regeneration & 30{,}000\,ms, self-rescheduling & \srccode{flux-domain-manager/\allowbreak{}src/\allowbreak{}services/\allowbreak{}domainService.js}{201-208} \\
Main pool minimum candidates & 10 (throws below this) & \srccode{flux-domain-manager/\allowbreak{}src/\allowbreak{}services/\allowbreak{}domainService.js}{181-183} \\
CDM certificate obtain/renew loop & 15\,min (900{,}000\,ms) & \srccode{flux-domain-manager/\allowbreak{}src/\allowbreak{}services/\allowbreak{}domainService.js}{1046-1088} \\
\bottomrule
\end{tabular}
\end{table}

Because the G-mode control-plane grace window (90\,s) is a much slower mechanism than HAProxy's
own data-plane health check ($\approx$6\,s detection floor), an already-open connection to a dead
sticky IP typically stops receiving traffic within HAProxy's own $\approx$6--9\,s window; the 90\,s
constant instead bounds how long \FDM keeps re-publishing the same, possibly-dead IP into the
configuration before attempting election of a new one
\srccode{flux-domain-manager/\allowbreak{}src/\allowbreak{}services/\allowbreak{}domainService.js}{253-321}. \shipped

\subsection{Certificates}
\label{sec:networking-certs}

Every certificate in this path is issued by Let's Encrypt, exclusively, via \texttt{certbot
certonly -\allowbreak{}-\allowbreak{}standalone -\allowbreak{}-\allowbreak{}http-\allowbreak{}01-\allowbreak{}port=8787}
\srccode{flux-domain-manager/\allowbreak{}src/\allowbreak{}services/\allowbreak{}domain/\allowbreak{}cert.js}{38-46}; no wildcard issuance is coded —
only exact-domain HTTP-01 challenges. The ACME account state (keys, order history) lives wherever
\texttt{certbot} itself keeps it, under \texttt{/\allowbreak{}etc/\allowbreak{}letsencrypt/\allowbreak{}}, outside \FDM's own code. The
artifact \texttt{HAProxy} actually serves is a single, world-readable-by-haproxy file per domain,
\texttt{/etc/ssl/\{certFolder\}/\{domain\}.pem}, formed by concatenating \texttt{fullchain.pem}
and \texttt{privkey.pem} — public certificate and private key in one file
\srccode{flux-domain-manager/\allowbreak{}src/\allowbreak{}services/\allowbreak{}domain/\allowbreak{}cert.js}{38-46}
\srccode{flux-domain-manager/\allowbreak{}src/\allowbreak{}services/\allowbreak{}haproxyTemplate.js}{141-150}. \shipped

Issuance is technically permissionless: any host that can bind port 8787 for the HTTP-01
challenge and run \texttt{sudo certbot} can obtain a certificate for a domain it is authoritative
for. Operationally, in the deployed fleet, only the 16 Flux-Foundation-run \FDM/CDM hosts ever do
so; an ordinary Cumulus, Nimbus or Stratus node operator never runs \FDM or its certificate-only
sibling \srccode{flux-fdm-infrastructure/\allowbreak{}ansible/\allowbreak{}inventory/\allowbreak{}hosts.yaml}{1-100}. Two PM2 processes
run from the same repository on the per-shard app hosts: \texttt{FDM} itself
(\texttt{manage\allowbreak{}Certificate\allowbreak{}Only:\allowbreak{} false}, handles HAProxy) and, on the EU-private-network hosts, a
second clone named \texttt{CDM} ("cert-domain-manager", \texttt{manage\allowbreak{}Certificate\allowbreak{}Only:\allowbreak{} true}),
dedicated purely to issuance and renewal
\srccode{flux-domain-manager/\allowbreak{}deployment/\allowbreak{}fdm\_setup.yml}{1-9,140-259}. CDM's loop runs every
15\,min, classifying each custom domain into skip/obtain/renew — renewing when fewer than 30 days
of validity remain — and deleting any certificate that expired more than 30 days ago
\srccode{flux-domain-manager/\allowbreak{}src/\allowbreak{}services/\allowbreak{}domain/\allowbreak{}cert.js}{18-36,91-126,214-249,308-340}. \shipped

\subsection{\fluxdnsd mesh resolution: a local file-tail resolver, not a gossip protocol}
\label{sec:networking-dnsd}

The name \fluxdnsd invites a reading in which nodes discover each other's application members by
talking to one another — a gossip or peer-to-peer resolution protocol. That reading is wrong, and
the paper states this plainly because the mistake is an easy one to make from the name alone.
\fluxdnsd is a purely local, single-node process. It holds an in-memory view built from exactly one
JSON file, \texttt{/\allowbreak{}var/\allowbreak{}lib/\allowbreak{}flux-mesh/\allowbreak{}resolver/\allowbreak{}membership.\allowbreak{}json}, and answers
\texttt{*.mesh.flux} queries from that view, scoped to the querying container's source IP
\srccode{flux-dnsd/\allowbreak{}src/\allowbreak{}snapshot.rs}{319-363,79-84}. It never opens a connection to another
\fluxdnsd instance, never queries a central server at request time, and implements no notion of
node-to-node consensus \srccode{flux-dnsd/\allowbreak{}src/\allowbreak{}snapshot.rs}{319-363}. The file is written by
\FluxOS, which computes the per-node membership view and installs it with a temporary-file-then-
rename, an atomic operation \fluxdnsd's file watcher depends on: only an inotify
\texttt{IN\_MOVED\_TO} event whose destination is \texttt{membership.\allowbreak{}json} triggers a reload —
a plain write-in-place is inert by construction
\srccode{flux-dnsd/\allowbreak{}src/\allowbreak{}watcher.rs}{20-69}. As the component's own documentation states,
"FluxOS's database is the authority and rewrites it on boot-sync"
\srcdoc{flux-dnsd/README.md:75-85}. \shipped

The cross-node propagation that makes membership on one node visible to a container on another —
the actual mechanism a reader might expect "mesh" to name — happens entirely inside \FluxOS,
outside \fluxdnsd. From \fluxdnsd's own source this propagation mechanism is opaque: the daemon
"only accepts or rejects [the file] whole," never repairing, merging, or arbitrating between
disagreeing sources \srcdoc{flux-dnsd/README.md:83-85}
\srccode{flux-dnsd/\allowbreak{}src/\allowbreak{}snapshot.rs}{72-77}. A snapshot whose generation number does not strictly
exceed the currently loaded view's is rejected outright, so a stale or replayed file cannot
regress the served view \srccode{flux-dnsd/\allowbreak{}src/\allowbreak{}snapshot.rs}{397-416}. Answers carry a fixed
5-second TTL, justified in the source as "membership changes must propagate and nothing in the
path caches anyway" \srccode{flux-dnsd/\allowbreak{}src/\allowbreak{}main.rs}{57-60}. If no valid snapshot has ever
arrived, or the source IP is unscoped, the daemon answers \texttt{REFUSED} rather than
\texttt{SERVFAIL}\slash\texttt{NXDOMAIN}, so a client's stub resolver fails over to its secondary
nameserver \srccode{flux-dnsd/\allowbreak{}src/\allowbreak{}server.rs}{82-90}. \shipped

\subsection{flux-dns-gateway: an mTLS proxy in front of the \PDNS key}
\label{sec:networking-gateway}

\texttt{flux-\allowbreak{}dns-\allowbreak{}gateway} is an mTLS-fronted proxy that lets the 16 \FDM instances write DNS
records into \PDNS without any of them ever holding the \PDNS API key themselves
\srcdoc{flux-dns-gateway/README.md:1-29}. Nginx terminates and verifies the client certificate and
sets \texttt{X-\allowbreak{}SSL-\allowbreak{}Client-\allowbreak{}Verify}\slash\texttt{X-\allowbreak{}SSL-\allowbreak{}Client-\allowbreak{}CN} headers that the FastAPI backend trusts
unconditionally, with no re-verification of the certificate itself in the Python layer
\srccode{flux-dns-gateway/\allowbreak{}src/\allowbreak{}dns\_gateway/\allowbreak{}api/\allowbreak{}middleware/\allowbreak{}auth.py}{18-69}. Each client's ACL is
enforced by \texttt{client\_cn}, \texttt{allowed\_zones}, \texttt{allowed\_record\_types}, and
\texttt{allowed\_prefixes}, deny-by-default for any unrecognised \texttt{client\_cn}
\srccode{flux-dns-gateway/\allowbreak{}src/\allowbreak{}dns\_gateway/\allowbreak{}config/\allowbreak{}acl.py}{13-97,99-151}. \shipped

The README's example ACL configuration also documents two further safety fields,
\texttt{max\_records\_per\_zone} and \texttt{rate\_limit\_per\_minute}
\srcdoc{flux-dns-gateway/README.md:259-274}. Neither has any corresponding reader anywhere in the
source: the \texttt{ACLConfig}\slash\texttt{FDMPermissions} parser reads only the four fields listed
above \srccode{flux-dns-gateway/\allowbreak{}src/\allowbreak{}dns\_gateway/\allowbreak{}config/\allowbreak{}acl.py}{21-30}. A documented per-zone
record cap and a documented per-client rate limit are, as deployed, no-ops. \shipped

\subsection{flux-geo-cert: a citation trap}
\label{sec:networking-geocert}

\texttt{flux-geo-cert}'s name and its corpus label ("Certificate Orchestrator for Flux Geo CDN")
suggest node-location attestation. It is not that. It is a Python service that automates Let's
Encrypt DNS-01 certificate renewal for the geo-routed CDN domain and pushes the resulting
certificates over SSH to a small set of CDN nodes (Germany, US, Hong Kong)
\srcdoc{flux-geo-cert/README.md:1-13}
\srccode{flux-fdm-infrastructure/\allowbreak{}ansible/\allowbreak{}inventory/\allowbreak{}hosts.yaml}{155-177}. It has no inbound network
interface at all — "pure outbound service" \srcdoc{flux-geo-cert/ARCHITECTURE.md:54} — and
performs no measurement, claim, or signature about where any Flux node physically is. No code
anywhere in the repositories surveyed for this section implements cryptographic attestation of
node geography; the geolocation machinery that actually feeds placement decisions
(fault-domain counting, residential-versus-hosting classification) lives in
\texttt{fluxos-\allowbreak{}network-\allowbreak{}policy}'s \texttt{iplocation.\allowbreak{}bin.\allowbreak{}gz} artifact
(\S\ref{sec:networking-policy}), a build-time table sourced from RIR delegated-extended files and
commercial geolocation data, not from any on-node measurement. This paper does not cite
\texttt{flux-geo-cert} for, and makes no claim anywhere that, node geography is cryptographically
attested.

\subsection{fluxos-network-policy: the governance-relevant lever}
\label{sec:networking-policy}

\texttt{fluxos-\allowbreak{}network-\allowbreak{}policy} is a small set of JSON/binary documents — blocked repositories,
vetted repositories, a tampering blocklist, an enterprise-node owner allow-list, and a
$\sim$2.0-million-row IP$\to$(organisation, country, region) geolocation table — that every
\FluxOS node fetches on a fixed schedule over plain HTTPS and enforces immediately, fleet-wide,
with no staged rollout \srcdoc{fluxos-network-policy/README.md:1-9,11-17,30-44}. Fetch intervals
are 6\,h for the blocked-repositories and enterprise-node documents and 12\,h for the tampering
blocklist, with the geolocation table refreshed on a daily conditional check
\srcdoc{fluxos-network-policy/README.md:34-35}
\srccode{flux/\allowbreak{}ZelBack/\allowbreak{}src/\allowbreak{}services/\allowbreak{}utils/\allowbreak{}cacheManager.js}{126-129}
\srccode{flux/\allowbreak{}ZelBack/\allowbreak{}src/\allowbreak{}services/\allowbreak{}utils/\allowbreak{}enterpriseConfig.js}{12}
\srccode{flux/\allowbreak{}ZelBack/\allowbreak{}src/\allowbreak{}services/\allowbreak{}appTamperingBlocklistService.js}{11}. A node that cannot reach the repository keeps
enforcing whatever it last successfully fetched, indefinitely and across restarts — failure here
is neither fail-open nor fail-closed-to-empty but fail-frozen, and a node that has never
successfully fetched anything declines to answer rather than guess
\srcdoc{fluxos-network-policy/README.md:33-44}. \shipped

The documents are unsigned and served over plain HTTPS; by the component's own account, "their
authenticity rests on GitHub and on who can merge to this repo"
\srcdoc{fluxos-network-policy/README.md:320-326}. Write access is restricted to GitHub
organisation owners, and there is deliberately no second-reviewer requirement: CI runs a shape
validator and a dual-write mirror check only, reaching green in roughly 15 seconds, after which
the author of the pull request may merge their own change
\srcdoc{fluxos-network-policy/README.md:297-303}
\srccode{fluxos-network-policy/\allowbreak{}.github/\allowbreak{}CODEOWNERS}{1-7}
\srccode{fluxos-network-policy/\allowbreak{}.github/\allowbreak{}workflows/\allowbreak{}validate.yml}{1-27}. The stated reasoning is
explicit: "the moment you most need to change policy is during an incident, and a rule requiring a
second person is one that gets bypassed at 3am"
\srcdoc{fluxos-network-policy/README.md:300-302}. Once merged, every \FluxOS node enforces
the new document on its own next scheduled fetch — within 7--13\,h depending on the document, the
fetch interval plus the hourly sweep (\S\ref{sec:governance}), with
no canary and no staged rollout \srcdoc{fluxos-network-policy/README.md:3-5,34-35}. \shipped

Concretely, a single merged pull request can: block or unblock any container image, namespace,
owner address, or 64-hex application hash network-wide (\texttt{blockedrepositories.\allowbreak{}json},
\texttt{vettedrepositories.\allowbreak{}json}); grant or revoke which owner addresses may install applications
on which enterprise-designated nodes (\texttt{enterprisenodes.\allowbreak{}json}); mark a node's collateral
transaction hash as blocked for tampering (\texttt{tamperingblockednodes.\allowbreak{}json}), a mechanism the
evidence describes as able to DOS a node's collateral for that offense; and, via the monthly
geolocation-baseline pull request, reclassify which organisations or regions count as residential
versus hosting for placement fault-domain counting and enforcement purposes
\srcdoc{fluxos-network-policy/README.md:320-326,99-172}. This is, by the component's own
documentation, "the sharpest single point of centralized, unsigned control" identified across the
network- and addressing-plane components covered in this section
\srcdoc{fluxos-network-policy/README.md:320-326}; \S\ref{sec:governance} argues directly from this
lever when it assesses censorship- and Sybil-resistance at the measured node distribution. \shipped

\subsection{The centralization inventory}
\label{sec:networking-inventory}

Table~\ref{tab:centralization} states, without further comment, which surfaces described in this
section are Foundation-operated, what depends on each, and what happens if each is unavailable.

\begin{table}[H]
\centering
\small
\renewcommand{\arraystretch}{1.2}
\caption{Centralization surfaces in the reachability and naming path.}
\label{tab:centralization}
\begin{tabular}{p{2.6cm}p{2.1cm}p{4.4cm}p{4.6cm}}
\toprule
Surface & Operator & Depends on it & If unavailable \\
\midrule
\texttt{api.\allowbreak{}runonflux.\allowbreak{}io} (application discovery) --- the \emph{configured default},
not a unique source & Flux Foundation operates this endpoint; the route it serves is on every node & Every \FDM/CDM
control-plane loop reads app specs, locations and permanent messages from one configured base URL. The
route is an ordinary \FluxOS{} one served by all \num{6489} nodes \srccode{flux/\allowbreak{}ZelBack/\allowbreak{}src/\allowbreak{}routes.js}{402}
and the host is a constructor parameter \srccode{flux-domain-manager/\allowbreak{}src/\allowbreak{}services/\allowbreak{}flux/\allowbreak{}dataFetcher.js}{886},
so an independent operator repoints it by configuration & \textbf{Resilience, not centralisation:} no
automatic failover to another node is coded, so a loop whose configured endpoint is unreachable keeps
polling, logs errors and serves the last-known HAProxy configuration (fail-static)
\srccode{flux-domain-manager/\allowbreak{}src/\allowbreak{}services/\allowbreak{}flux/\allowbreak{}dataFetcher.js}{47-91} \\
\addlinespace
Foundation-private \texttt{10.100.0.0/24} (WireGuard hub, \PDNS masters, DNS Gateway, Arcane load
balancers, SAS decrypt endpoint) & Flux Foundation & All non-EU \FDM/CDM and
\texttt{flux-\allowbreak{}apps-\allowbreak{}dns-\allowbreak{}manager} instances, reachable outside the EU segment only through a single
WireGuard hub, \texttt{wireguard-\allowbreak{}eu-\allowbreak{}01} (\texttt{10.100.0.167}, public
\texttt{5.39.57.49:\allowbreak{}51000}) \srccode{flux-fdm-infrastructure/\allowbreak{}ansible/\allowbreak{}inventory/\allowbreak{}hosts.yaml}{180-185}
& Custom-domain certificate issuance, enterprise application-spec decryption, and game-application
DNS publication all stop for every host that reaches this segment only through that hub
\srccode{flux-domain-manager/\allowbreak{}src/\allowbreak{}services/\allowbreak{}domainService.js}{1099-1105} \\
\addlinespace
16-host \FDM/CDM fleet (6 EU direct-private, 6 US and 4 Singapore over WireGuard)
\srccode{flux-fdm-infrastructure/\allowbreak{}ansible/\allowbreak{}inventory/\allowbreak{}hosts.yaml}{1-100} & Flux Foundation & HAProxy
routing and TLS termination for every \texttt{app(2).\allowbreak{}runonflux.\allowbreak{}io} and custom-domain hostname
network-wide & Stale-but-serving: a failed configuration validation leaves the previous
configuration running, but no new application is routed or certified until a host recovers
\srccode{flux-domain-manager/\allowbreak{}src/\allowbreak{}services/\allowbreak{}domainService.js}{196-200} \\
\addlinespace
\texttt{flux-\allowbreak{}apps-\allowbreak{}dns-\allowbreak{}manager} (single, unreplicated instance on \texttt{flux-cert-1} =
\texttt{10.100.0.172}) & Flux Foundation & Direct node-to-node DNS for $\sim$26 G-mode
game-server application types
\srccode{flux-apps-dns-manager/\allowbreak{}config/\allowbreak{}gamesConfig.js}{30-57} & No coded high availability,
leader election, or standby process; game-application \texttt{A} records stop updating until the
single host/co-located DNS Gateway returns
\srccode{flux-apps-dns-manager/\allowbreak{}ansible/\allowbreak{}inventory.yaml}{6-7} \\
\addlinespace
Let's Encrypt (ACME certificate authority) & External (ISRG) & Every TLS-terminated
\texttt{app(2).\allowbreak{}runonflux.\allowbreak{}io} and custom-domain certificate; exclusively used, no alternate CA
coded \srccode{flux-domain-manager/\allowbreak{}src/\allowbreak{}services/\allowbreak{}domain/\allowbreak{}cert.js}{38-46} & No fallback CA; the CDM
renewal loop retries every 15\,min indefinitely, and the 30-day renewal floor provides slack, but
new domains cannot be certified at all during an outage \\
\addlinespace
Cloudflare (where \texttt{config.\allowbreak{}cloudflare.\allowbreak{}enabled}) & External (Cloudflare) & DNS resolution for
domains on the Cloudflare-selected path, used strictly DNS-only (grey-cloud), never as a proxy or
CDN \srccode{flux-domain-manager/\allowbreak{}src/\allowbreak{}services/\allowbreak{}domain/\allowbreak{}cert.js}{170-181} & No fallback registrar or
DNS backend is coded per domain; whichever hostnames are Cloudflare-hosted stop resolving \\
\bottomrule
\end{tabular}
\end{table}

\subsection{Node network configuration}
\label{sec:networking-nodeconfig}

At install and reconfiguration time, network interfaces are discovered automatically: a
\texttt{probert}-backed (netlink and udev) scan produces interface objects carrying type, name,
state, index, MAC, vendor, model, address, VLAN and bridge identifiers with no operator input
\srccode{fluxos-iso-networking-shared/\allowbreak{}src/\allowbreak{}flux\_networking\_shared/\allowbreak{}models/\allowbreak{}network.py}{326-388}.
What the operator configures, per interface, is the allocation mode — DHCP or static — through a
TUI screen; choosing static reveals a form for subnet, address, gateway and DNS that is
blur-validated before the change may be saved
\srccode{fluxos-iso-networking/\allowbreak{}src/\allowbreak{}flux\_iso\_networking/\allowbreak{}screens/\allowbreak{}interface\_modal\_screen.py}{54,79-84,200-284}.
Per-interface traffic shaping is likewise operator-set, from a closed set of rates
($\{35,75,135,250\}$\,Mbps), written to a policy file and applied by a separate, privileged D-Bus
helper because the configuring process itself runs without the capability to apply it directly
\srccode{fluxos-iso-networking-shared/\allowbreak{}src/\allowbreak{}flux\_networking\_shared/\allowbreak{}models/\allowbreak{}network.py}{77,83,139-161,224}.
\shipped

The stack is IPv4-only by construction, not merely by current deployment. Every network dataclass
— \texttt{Route}, \texttt{SourceAddress}, \texttt{NetworkInterface.\allowbreak{}address} — is typed
\texttt{IPv4Network}\slash\texttt{IPv4Address}\slash\texttt{IPv4Interface} and nothing else, and the shared
connectivity helper pins \texttt{TCPConnector(family=\allowbreak{}AF\_INET)} explicitly
\srccode{fluxos-iso-networking-shared/\allowbreak{}src/\allowbreak{}flux\_networking\_shared/\allowbreak{}models/\allowbreak{}network.py}{17-29,326-337}
\srccode{fluxos-iso-networking-shared/\allowbreak{}src/\allowbreak{}flux\_networking\_shared/\allowbreak{}helpers.py}{147,242}. The
measured fleet crawl underlying the network-wide baseline in \S\ref{sec:baseline} reports
2{,}322 IPv4 addresses and zero IPv6 or onion (Tor hidden-service) addresses
\srcmeasured{api.runonflux.io/\allowbreak{}daemon/\allowbreak{}getzelnodecount}{2026-09-03}. \shipped

A reachability design for NAT'd and non-public nodes exists only as a docstring. The helper
\texttt{is\_directly\_public()} states that a node with a globally-routable public IP bound to one
of its own interfaces would act as a "Tor SOCKS-only hub," while a NAT'd node would instead run a
"Tor hidden service" to become reachable
\srccode{fluxos-iso-networking-shared/\allowbreak{}src/\allowbreak{}flux\_networking\_shared/\allowbreak{}models/\allowbreak{}network.py}{461-479}.
No Tor-related code exists anywhere else in either the installer or the shared networking library;
the design is fully specified in prose and entirely unimplemented. \vision

\subsection{Private deployments (\VPON)}
\label{sec:networking-vpon}

The roadmap lists private deployments, referred to as \VPON s (private overlay networks), as a
Q4 2026 item, described as "the complement to private payments and chain privacy"
\srcroadmap{Private deployments (VPONs), Q4 2026}; the roadmap's own timeline artifact labels the
same item "Private overlay networks," sub-labelled "VPONs $\cdot$ private deployments"
\srcdoc{flux-roadmap/public/timeline.html:172-173}. No implementation of a \VPON was found in any
repository surveyed for this section or flagged elsewhere in the corpus; the roadmap evidence
itself notes no repository name observed anywhere in the codebase corresponds to a \VPON
\srcdoc{flux-roadmap/public/flux-roadmap-public.md:83}. \planned

Because no such system exists to describe, this section states what a \VPON would need to provide
rather than how one works. At minimum, consistent with the notation this paper uses for the term
and with the roadmap's framing of \VPON s as private overlays rather than a naming variant on the
public network:
\begin{itemize}
  \item \textbf{Membership by key, not by discoverability.} The reachability path in
  \S\ref{sec:networking-path} and the mesh namespace in \S\ref{sec:networking-dnsd} both resolve
  any application's members to any container scoped correctly on the querying node; a \VPON would
  need membership gated by possession of a private key rather than resolvable by any party who can
  reach the mesh namespace or the public API \srcinferred.
  \item \textbf{Reachability-plane confidentiality.} If a \VPON exists to complement private
  payment (\S\ref{sec:private-payment}) and chain privacy, then the infrastructure this section
  describes — the discovery API, \FDM, and \PDNS, all Foundation-operated — should not be
  positioned to learn which nodes host a private deployment or on whose behalf, since today's
  reachability path is Foundation-legible end to end by the construction in
  \S\ref{sec:networking-path} \srcinferred.
\end{itemize}
The roadmap sequences \VPON s alongside private payments and the chain-privacy work without
stating an explicit blocking dependency between them
\srcdoc{flux-roadmap/public/flux-roadmap-public.md:81,84}. Beyond this framing, no wire protocol,
key-management scheme, or placement-isolation mechanism for a \VPON is specified anywhere in the
roadmap or in the repositories surveyed. The construction is open.

What does ship, today, are WireGuard-based transports inside \CumulusVPN — a single VPN product,
not a general per-application private-overlay primitive — that a future \VPON design could draw
on as building blocks:
\begin{itemize}
  \item \textbf{AmneziaWG obfuscation.} A fork of \texttt{wireguard-go} that reshapes WireGuard's
  handshake framing without altering its cryptography, wired up as an additive listener sharing
  the vanilla gateway's server key. Gated off by default via an environment flag, but the fleet
  deployment manifest sets it on by default
  \srccode{cumulusvpn/\allowbreak{}gateway/\allowbreak{}internal/\allowbreak{}wg/\allowbreak{}device.go}{19,67-124}
  \srccode{cumulusvpn/\allowbreak{}deploy/\allowbreak{}countries.yaml}{23}. \shipped, disabled by default in code, enabled
  by default in the deployed fleet manifest.
  \item \textbf{WireGuard-over-TLS ("stealth" transport, port 443).} A TLS session terminated on a
  TCP port and bridged, byte-for-byte, to the local WireGuard UDP listener; the client does not
  verify the certificate chain, and trust is anchored entirely in the inner WireGuard handshake
  \srccode{cumulusvpn/\allowbreak{}gateway/\allowbreak{}internal/\allowbreak{}tlsrelay/\allowbreak{}relay.go}{1-16,229-236}. Also gated off by default
  in code, also on by default in the fleet manifest
  \srccode{cumulusvpn/\allowbreak{}deploy/\allowbreak{}countries.yaml}{24}. \shipped, same disabled-in-code/enabled-in-fleet
  pattern as above.
  \item \textbf{Multi-hop.} Implemented client-side only, with no corresponding change to the
  gateway protocol: the client runs two stacked WireGuard interfaces under the same key, and the
  entry hop decrypts only the outer layer before forwarding opaque ciphertext to the exit hop
  \srcdoc{cumulusvpn/docs/11-multihop.md:8-33}. Because both hops see the same tunnel key in this
  version, an adversary controlling both a user's chosen hops could correlate traffic by key — a
  gap the design documentation states and leaves open, not a hypothetical one
  \srcdoc{cumulusvpn/docs/11-multihop.md:47-58}. \shipped, client-side, with a stated unmitigated
  correlation risk.
\end{itemize}
Nothing in the evidence surveyed shows any of these three transports wired into a general,
per-application \VPON primitive; they are cited here only as the concrete pieces that exist and
that a future \VPON specification would not need to invent from nothing.

%% Drain this section's float queue before the next begins.
\FloatBarrier
